% =====================================================================
% RAOD++ / iso-cost study — Discussion
% Claims here are limited to what Section 4 measured.
% =====================================================================

\section{Discussion}
\label{sec:discussion}

\subsection{Why the two axes disagree}
\label{sec:discussion:mechanism}

The central result is that the same question has two answers depending on the resource
being spent. The power measurements suggest why. Raising the input from $640$ to $1536$
pixels multiplies the pixel count by $5.8$ but moves board power from $57.3$ to
$108.4$\,W, a factor of $1.9$; moving from the nano to the small model at $1280$ pixels
moves it from $88.0$ to $146.4$\,W while also increasing latency. Board power includes an
idle floor of $37.3$\,W on this device, and removing it sharpens the contrast rather than
softening it: the marginal draw of inference rises from $19.9$ to $71.1$\,W along the
resolution arm and from $50.7$ to $109.1$\,W along the capacity arm. Reusing one set of weights over
more pixels is arithmetically heavier but memory-cheaper than moving a larger set of
weights and activations, and on modern accelerators energy tracks data movement more
closely than it tracks arithmetic.

Latency at batch one does not expose this difference, and for a specific reason: a model of
this size does not saturate the device. Raising the input from $640$ to $1280$ multiplies
pixels by four while latency rises only from $4.98$ to $6.50$\,ms, a factor of $1.31$,
because a substantial share of the wall-clock time is spent in preprocessing, non-maximum
suppression and framework overhead rather than in convolutions. Under those conditions a
larger model is nearly free in wall-clock terms and expensive in joules, which is precisely
the ordering Table~\ref{tab:isocost} reports.

This also bounds the claim. On hardware where compute is the binding constraint --- an
embedded accelerator, a larger backbone, or batched inference --- the latency ranking could
change, while the energy argument, resting on data movement, is the more likely of the two
to transfer. We measured neither case and do not claim either.

\subsection{Which energy a deployed system actually spends}
\label{sec:discussion:energy}

Total board energy, $\bar{P} L$, charges a configuration for every watt the card draws
while it works. That is the right quantity only if the detector occupies the whole frame
period. An airborne platform with a 30\,Hz camera has a frame period of $33.3$\,ms, and the
detectors measured here finish in $5$ to $7$\,ms, so for most of each period the device
sits at its idle floor whichever configuration is installed. Writing $T$ for the frame
period and $t$ for the time the detector runs, the energy spent per frame is closer to
\begin{equation}
  E_{\text{frame}} \;=\; P_{\text{idle}} T \;+\; \big(P_{\text{active}} - P_{\text{idle}}\big)\, t ,
  \label{eq:frame-energy}
\end{equation}
in which the first term is the same for every configuration and only the second
discriminates between them. That second term is exactly the marginal energy reported as a
robustness check in Section~\ref{sec:results:isocost}.

This does not change the verdict, since the ordering is the same under both definitions,
and the pre-registered comparison was fixed on total board energy before the idle floor was
measured. It does change which number a practitioner should read. At the duty cycles a UAV
detector runs at, the marginal figure is the more decisive one, and its margin is the wider
of the two: resolution leads capacity by $0.0457$ above the idle floor against $0.0331$ on
total power. Where a detector saturates the frame period instead, the two converge and
total energy becomes the right measure. Reporting both, and stating the duty cycle a claim
assumes, costs nothing.

\subsection{What the field's reporting conventions hide}
\label{sec:discussion:conventions}

Published work on small-object detection from UAVs reports gains of roughly $+2.8$ to
$+10.5$ mAP$_{50}$ from architectural modifications. Raising the input resolution of an
unmodified nano detector, with no change to the architecture and no additional data, yields
$+0.1386$ mAP$_{50\text{-}95}$ over the $640$ baseline, a relative gain of $74\%$. The two
numbers are not directly comparable, which is exactly the point: resolution is rarely
controlled in those comparisons, so an architectural gain measured at one input size and a
baseline trained at another are reported in the same table.

Two conventions contribute. First, parameters and FLOPs are reported as proxies for
deployment cost, and in the batch-one regime they predict latency poorly: the P2 variant is
slower than the small model at the same input size despite having fewer parameters, while
drawing less power and therefore winning on energy where it loses on latency, and
resolution changes move FLOPs far more than they move measured latency. Second, energy is
almost never reported at all, although for an airborne platform the binding budget is
frequently the battery rather than the frame period.

We are not arguing that architectural work is misdirected. We are arguing that a comparison
that does not fix a measured cost cannot tell the reader which allocation was responsible
for the gain, and that reporting both axes is cheap relative to the training that produced
the numbers.

For the same reason this paper reports no comparison against published detectors on these
benchmarks, and the omission is deliberate. A fair comparison would require retraining
those methods at matched input size on the same split, and measuring their latency and
energy on this machine; published numbers were obtained at input sizes the authors chose,
on hardware we do not have, and with cost reported as parameters or FLOPs where it is
reported at all. Placing them in a table beside ours would reproduce exactly the confound
this work is about. What we claim is a controlled statement about where a fixed budget
should be spent within one detector family, not a position on a leaderboard, and the two
require different experiments.

\subsection{Three negative results, and why they are load-bearing}
\label{sec:discussion:negatives}

The resolution finding is more informative because of what preceded it.

Enriching the training data with retrieved, semantically matched instances did not help.
All five conditions of the control matrix land within $0.0025$ mAP$_{50\text{-}95}$ of the
baseline and all fall below it; the retrieval condition beats the no-data control by
$0.0010$, statistically detectable and practically negligible. Raising the intervention from
$2.7\%$ to $9.5\%$ of the existing boxes leaves accuracy flat and then degrades it, so the
result is not an artefact of scale. The random-paste control shares the placement engine,
budget and filters with the retrieval condition and differs only in how the neighbour is
chosen, so the null result concerns the semantics of the retrieval and not the mechanics of
the pasting.

The retrieval stage itself is not inert. Its information content is inversely related to a
class's share of the gallery, from $23.4$ times chance for the rarest class to exactly
chance for the most frequent. That pattern is a useful design constraint for
retrieval-augmented detection in general, and it also explains part of the null result:
the classes for which retrieval carries information are the ones for which the gallery
supplies fewest examples to paste.

The second negative result concerns confidence. If small objects were merely under-ranked,
a calibration correction would recover them at no cost. Measured over ten seeds, the
size-conditioned confidence gap spans only $0.0295$ and runs in the direction opposite to
that hypothesis. The detector's low confidence on small objects is an accurate report, not a
bias. What the same measurement does show is severe: at a deployment threshold of $0.5$,
recall on the smallest size band is $0.053$. A detector reported by its mAP alone can miss
most of the objects it is deployed to find, and this gap between integrated and operating
point behaviour deserves to be reported routinely.

The third is our own first hypothesis. We initially predicted that resolution would beat
capacity at fixed latency, and that prediction was rejected. The error was procedural: the
comparison pairs were fixed before the cost axis was measured, and clean measurement showed
that two conditions intended as a matched pair differed by $18\%$ in latency. Replacing
pairwise comparison with the curve-based procedure of Section~\ref{sec:method:isocost} is
what made the question answerable.

\subsection{Practical guidance}
\label{sec:discussion:guidance}

For a real-time UAV detector of this class, the results support three concrete
recommendations. When the binding constraint is the frame period, spend the budget on model
capacity. When it is the energy budget, spend it on input resolution. In both cases, prefer
either to adding a P2 detection head, which lost on both axes here at a comparable cost.

Two caveats attach. The accuracy--resolution curve had not saturated at the largest input
tested, so the ceiling was not located; the recommendation is about allocation under a
budget, not an argument to raise resolution without limit. And the advantage of one arm over
the other narrows as the budget grows, from $-0.0314$ to $-0.0101$ on latency and from
$+0.0502$ to $+0.0229$ on energy across the measured intervals, so the choice matters most
where budgets are tightest.

\subsection{Limitations}
\label{sec:discussion:limitations}

Cost was measured on one desktop GPU at batch one. That is the deployment mode of interest,
but the device is not saturated by a model of this size, and the latency ordering may differ
on hardware where compute binds. No embedded validation was performed and none is claimed. Energy is NVML board power rather
than system wall power.

Two training conditions were not held constant and should temper how cleanly the factors
are separated. The physical batch falls from 64 at 640 pixels to 6 at 1536, so
batch-normalisation statistics are estimated over samples of different size along the
resolution arm, and along the capacity arm as well. And early stopping reads the same split
the results are reported on. Neither can favour one arm over the other, since both apply
under one rule to every condition, but both mean the arms differ in slightly more than the
single factor each is named for.

The grid uses three seeds per condition; only the $640$ baseline was repeated over ten. The
seed-to-seed standard deviation is small, which is why differences of this magnitude are
resolvable, but confidence intervals should be read with the sample size in mind. All
experiments use a single detector family; whether the ordering holds for two-stage or
transformer detectors is untested.

The arms are compared only over the cost interval where both were measured. Those intervals
are finite, and the reported difference is an average over them rather than a statement
about every possible budget. Extending an arm changes the overlap, which is why the
confirmation runs of Section~\ref{sec:results:isocost} were added before any claim was made.

Finally, the finding that the answer depends on the scarce resource emerged after the data
were seen. It was restated as a hypothesis with a fixed success criterion before the
confirming experiments finished, and it is reported as confirmed on VisDrone and not confirmed on UAVDT, where both
directions replicated at comparable magnitude but neither reached corrected significance.
Two datasets in one domain remain two datasets, and generalisation beyond aerial imagery is
untested.
